\scanpage{152}
\mspage{101}

j'ai passé des heures avant d'en trouver un et d'y voir un peu clair ! Il suffit en somme d'exhiber un exemple d'une \unsure{flèches|flèche} dans $\mathrm{Cat}$

\[
 f:X\longrightarrow Y
\]

telle que

\[
\begin{aligned}
\text{a)}\quad &\underline{\unsure{\pi|\Pi}}_0(f)\overset{\mathrm{def}}{=}\underline{\pi}_0(X/\unsure{Y|y})\xrightarrow{\sim}Y\\
&\text{i.e. les }\pi_0({}_{y\backslash}X)\text{ sont 0-conn.}\\
&\hspace{7em}(\text{pour }y\in\mathrm{ob}\, Y)
\end{aligned}
\]

et telle que de plus

\[
\text{b)}\qquad X\text{ est }\unsure{W|\mathcal W}\text{-asph, et }Y\text{ non.}
\]

% vertical rule beside the argument through the next display
\noindent\makebox[0pt][r]{\rule[-4.8\baselineskip]{0.4pt}{5.6\baselineskip}\hspace{.65em}}%
La difficulté est la suivante, quand on cherche des contre-exemples ``idiots'' : c'est que l'hypothèse a) \unsure{Implique|implique} déjà, quand $X$ est $W_\infty$-asph. (ou seulement 1-connexe), que $Y$ est 1-connexe : le foncteur

\[
Y'\longrightarrow X'=X\times_Y Y'
\]

de la cat. des rev. étales $Y'$ de $Y$ vers celles des rev. étales $X'$ de $X$ est pleinement fidèle\unsure{$^{x}$} (vérification immédiate), d'où le résultat annoncé. Donc un contrexemple doit faire appel à l'existence d'un $\pi_i(Y,y)$ non nul, avec $i\geq2$. Cela \unsure{explique|exclut donc|exclut}
% archive sheet number: 152
